%-----------------------------------------------------------------------------------------------------------------------------------------------%
%	The MIT License (MIT)
%
%	Copyright (c) 2021 Jitin Nair
%
%	Permission is hereby granted, free of charge, to any person obtaining a copy
%	of this software and associated documentation files (the "Software"), to deal
%	in the Software without restriction, including without limitation the rights
%	to use, copy, modify, merge, publish, distribute, sublicense, and/or sell
%	copies of the Software, and to permit persons to whom the Software is
%	furnished to do so, subject to the following conditions:
%
%	THE SOFTWARE IS PROVIDED "AS IS", WITHOUT WARRANTY OF ANY KIND, EXPRESS OR
%	IMPLIED, INCLUDING BUT NOT LIMITED TO THE WARRANTIES OF MERCHANTABILITY,
%	FITNESS FOR A PARTICULAR PURPOSE AND NONINFRINGEMENT. IN NO EVENT SHALL THE
%	AUTHORS OR COPYRIGHT HOLDERS BE LIABLE FOR ANY CLAIM, DAMAGES OR OTHER
%	LIABILITY, WHETHER IN AN ACTION OF CONTRACT, TORT OR OTHERWISE, ARISING FROM,
%	OUT OF OR IN CONNECTION WITH THE SOFTWARE OR THE USE OR OTHER DEALINGS IN
%	THE SOFTWARE.
%
%
%-----------------------------------------------------------------------------------------------------------------------------------------------%

%----------------------------------------------------------------------------------------
%	DOCUMENT DEFINITION
%----------------------------------------------------------------------------------------

% article class because we want to fully customize the page and not use a cv template
\documentclass[a4paper,12pt]{article}

%----------------------------------------------------------------------------------------
%	PACKAGES
%----------------------------------------------------------------------------------------
\usepackage{url}
\usepackage{parskip}

%other packages for formatting
\RequirePackage{color}
\RequirePackage{graphicx}
\usepackage[usenames,dvipsnames]{xcolor}
\usepackage[scale=0.9]{geometry}

%tabularx environment
\usepackage{tabularx}

%for lists within experience section
\usepackage{enumitem}

% centered version of 'X' col. type
\newcolumntype{C}{>{\centering\arraybackslash}X}

%to prevent spillover of tabular into next pages
\usepackage{supertabular}
\usepackage{tabularx}
\newlength{\fullcollw}
\setlength{\fullcollw}{0.47\textwidth}

%custom \section
\usepackage{titlesec}
\usepackage{multicol}
\usepackage{multirow}

%CV Sections inspired by:
%http://stefano.italians.nl/archives/26
\definecolor{sectioncolor}{RGB}{0, 102, 204}  % Define blue color for sections
\titleformat{\section}{\Large\scshape\raggedright\color{sectioncolor}}{}{0em}{}[\titlerule]
\titlespacing{\section}{0pt}{10pt}{10pt}

%for publications
\usepackage[style=authoryear,sorting=ynt, maxbibnames=2]{biblatex}

%Setup hyperref package, and colours for links
\usepackage[unicode, draft=false]{hyperref}
\definecolor{linkcolour}{rgb}{0,0.2,0.6}
\hypersetup{colorlinks,breaklinks,urlcolor=linkcolour,linkcolor=linkcolour}
\addbibresource{citations.bib}
\setlength\bibitemsep{1em}

%for social icons
\usepackage{fontawesome5}

%debug page outer frames
%\usepackage{showframe}


% job listing environments
\newenvironment{jobshort}[2]
    {
    \begin{tabularx}{\linewidth}{@{}l X r@{}}
    \textbf{#1} & \hfill &  #2 \\[3.75pt]
    \end{tabularx}
    }
    {
    }

\newenvironment{joblong}[2]
    {
    \begin{tabularx}{\linewidth}{@{}l X r@{}}
    \textbf{#1} & \hfill &  #2 \\[3.75pt]
    \end{tabularx}
    \begin{minipage}[t]{\linewidth}
    \begin{itemize}[nosep,after=\strut, leftmargin=1em, itemsep=3pt,label=--]
    }
    {
    \end{itemize}
    \end{minipage}
    }



%----------------------------------------------------------------------------------------
%	BEGIN DOCUMENT
%----------------------------------------------------------------------------------------
\begin{document}

% non-numbered pages
\pagestyle{empty}

%----------------------------------------------------------------------------------------
% TITLE
%----------------------------------------------------------------------------------------

\begin{tabularx}{\linewidth}{@{} X r @{}}
\Huge{Hengbin Gao} & \multirow{7}{*}{\includegraphics[width=3cm]{profile.jpg}} \\[3pt]
\normalsize{\textit{Undergraduate}} & \\
\normalsize{\textit{Zhejiang University, Information Engineering}} & \\[10pt]
\href{mailto:Particle_Nebula@outlook.com}{\raisebox{-0.05\height}\faEnvelope \ Particle\_Nebula@outlook.com} \ $|$ \
\href{tel:+8619883109386}{\raisebox{-0.05\height}\faMobile \ (+86)\ 19883109386} & \\[3pt]
\href{mailto:3230105132@zju.edu.cn}{\raisebox{-0.05\height}\faEnvelope \ 3230105132@zju.edu.cn} \ $|$ \
\href{https://Particlela.github.io}{\raisebox{-0.05\height}\faGlobe \ Particlela.github.io} & \\[3pt]
\href{https://github.com/Particlela}{\raisebox{-0.05\height}\faGithub\ Particlela} & \\
\end{tabularx}
\vspace{10pt}

%----------------------------------------------------------------------------------------
% EXPERIENCE SECTIONS
%----------------------------------------------------------------------------------------

\section{Summary}
I am currently a third-year undergraduate student at the College of Information Science and Electronic Engineering, Zhejiang University, majoring in Information Engineering. As a member of the RoboMaster electrical control team, my research focuses on the control of series-elastic leg robots. I have also studied the control of four-switch buck-boost buffer capacitors. Moving forward, I am interested in exploring machine learning, digital signal/image processing, and embedded systems. Please feel free to contact me!

%----------------------------------------------------------------------------------------
%	EDUCATION
%----------------------------------------------------------------------------------------
\section{Education}
\begin{tabularx}{\linewidth}{@{}l X@{}}
2023 - 2027 & B.Eng. in \textbf{Information Engineering}, Zhejiang University \hfill \normalsize \\
& College of Information Science and Electronic Engineering \\
& Major GPA: 4.48/5.0 \\
& Overall GPA: 4.46/5.0 \\
\end{tabularx}

%----------------------------------------------------------------------------------------
%	HONORS AND AWARDS
%----------------------------------------------------------------------------------------
\section{Honors \& Awards}
\begin{tabularx}{\linewidth}{@{}l X@{}}
2026 & \textbf{Third Prize}, RoboMaster University Championship --- Infantry Robot Competition \\[3.75pt]
2026 & \textbf{Second Prize}, RoboMaster University Championship --- Sentry Robot Competition \\[3.75pt]
2025 & \textbf{First Prize}, RoboMaster University Championship --- National Final \textbf{(National Top 8)} \\[3.75pt]
2025 & \textbf{First Prize}, RoboMaster University Championship --- Infantry Robot Competition \\[3.75pt]
2025 & \textbf{Third Prize Scholarship}, National Talent Training Base \\[3.75pt]
2025 & \textbf{Third Prize Scholarship}, Zhejiang University \\[3.75pt]
2024 & \textbf{Provincial Government Scholarship}, Zhejiang University \\[3.75pt]
2024 & \textbf{Second Prize Scholarship}, Zhejiang University \\
\end{tabularx}

%Projects
\section{Projects}

\begin{tabularx}{\linewidth}{ @{}l r@{} }
\textbf{RoboMaster: Series-Elastic Leg Balancing Infantry Robot Control System} & \hfill \href{https://Particlela.github.io/project/}{Oct 2024 - Present} \\[3.75pt]
\multicolumn{2}{@{}X@{}}{Built on the STM32H7 platform for the RoboMaster series-elastic leg wheeled balancing infantry robot, based on a five-bar linkage leg mechanism and designed for high-dynamic locomotion. At the core control layer, a \textbf{variable-gain LQR balance control strategy} is adopted: a 10-state LQR controller is designed with its gain matrix and feedforward terms expressed as bivariate cubic polynomial functions of the left and right virtual leg lengths, enabling online real-time interpolation to ensure stable balancing across varying leg configurations. At the decision and behavior scheduling level, the system runs a \textbf{priority-based hierarchical state machine} covering Dead, Recovery, Mature, Flight, Jump, Onestep, Twostep, and Creep modes, supporting orderly switching and execution of complex maneuvers such as obstacle jumping and single/double-step climbing. For energy, the system integrates a supercapacitor and a buffer capacitor to form a two-stage energy buffer, and introduces a \textbf{quadratic-model-based power limiter} that dynamically scales velocity and yaw-axis angular velocity commands based on the real-time power limit from the referee system and remaining energy status, balancing maneuverability and system safety within a limited power budget, and ultimately forming a complete control loop integrating balance control, behavior management, and energy scheduling.}  \\
\end{tabularx}

\begin{tabularx}{\linewidth}{ @{}l r@{} }
\textbf{RoboMaster: Series-Elastic Leg Sentry Robot Control System} & \hfill \href{https://Particlela.github.io/project/}{Aug 2025 - Present} \\[3.75pt]
\multicolumn{2}{@{}X@{}}{Built upon the infantry balancing robot framework, the system is further upgraded to an autonomous navigation-oriented sentry robot, sharing the same series-elastic leg hardware platform and low-level control architecture, but shifting the control focus from manual operation to autonomous decision-making and high-level battlefield maneuvering. To achieve unmanned operation, the system adds a \textbf{chassis navigation command interface} for receiving motion reference commands from the upper-level NUC, enabling the sentry to completely eliminate dependence on remote control or keyboard/mouse input and independently execute navigation tasks. The sentry robot not only inherits the infantry's high-dynamic balancing capability but also possesses complete intelligence for autonomous planning and execution of complex maneuvers in unknown environments.}  \\
\end{tabularx}

\begin{tabularx}{\linewidth}{ @{}l r@{} }
\textbf{Bidirectional Four-Switch Buck-Boost Buffer Capacitor Control System (FSBB)} & \hfill \href{https://Particlela.github.io/project/}{June 2026 - Present} \\[3.75pt]
\multicolumn{2}{@{}X@{}}{Based on the STM32G4 MCU, this system designs and implements a high-bandwidth buffer-capacitor energy management controller for RoboMaster mobile robots. The main power stage adopts a bidirectional four-switch Buck-Boost converter topology, capable of dynamically absorbing and releasing energy to effectively suppress instantaneous bus voltage drops during high-dynamic maneuvers. At the control level, the system runs a \textbf{state-machine-based multi-loop control strategy}, integrating voltage-ratio feedforward compensation with three independent PI loops (constant-power charging, capacitor voltage regulation, and bus voltage support), enabling the controller to smoothly switch between charging, voltage regulation, and discharging modes based on bus conditions. Multi-channel ADC and DMA are used for high-speed acquisition of input voltage, capacitor voltage, input current, output current, and capacitor current, with precision calibration and complementary filtering to provide high-accuracy, low-latency real-time state feedback for each closed loop. Safety mechanisms integrate power ramp-up, over-voltage / under-voltage / over-current threshold protection, and automatic shutdown / recovery logic on CAN communication timeout or abnormal operating conditions, ensuring reliable hardware operation in complex electrical environments. Achieved key performance specifications: capacitor voltage operating range 11 V--28 V, maximum charging power 190 W, maximum output / supplement power 1200 W, capacitor current limit 70 A, input voltage range 15 V--28 V.}  \\
\end{tabularx}

\begin{tabularx}{\linewidth}{ @{}l r@{} }
\textbf{HW-Components: Universal Robot Control Library} & \hfill \href{https://Particlela.github.io/project/}{Oct 2024 - Present} \\[3.75pt]
\multicolumn{2}{@{}X@{}}{As a core contributor, I architected and maintained the embedded communication layer of HW-Components, a reusable robot control library targeting RoboMaster and general mobile robotic platforms.}  \\
\end{tabularx}

\begin{tabularx}{\linewidth}{@{}X@{}}
\begin{itemize}[nosep,after=\strut, leftmargin=1em, itemsep=3pt,label=--]
\item \textbf{Multi-Protocol Ranging Sensor Communication}: Implemented unidirectional and bidirectional data exchange between the STM32 MCU and various ranging modules (ultrasonic, ToF, infrared, etc.), providing flexible perception interfaces for different robotic configurations.
\item \textbf{Referee System Serial Link}: Maintained and optimized the UART-based communication protocol with the RoboMaster referee system, ensuring reliable reception of match-critical data such as robot status, hit points, firing permissions, and power limits.
\item \textbf{Energy Management over CAN}: Developed bidirectional CAN communication with the supercapacitor and buffer-capacitor subsystems, enabling real-time power monitoring, energy buffering control, and dynamic power allocation.
\item \textbf{Lever-Arm Compensation}: Implemented lever-arm compensation to correct kinematic errors caused by the physical offset between the IMU mounting position and the robot's center of rotation, improving state-estimation accuracy during rapid attitude changes.
\end{itemize}
\end{tabularx}

%----------------------------------------------------------------------------------------
%	PUBLICATIONS
%----------------------------------------------------------------------------------------
% \section{Publications}
% \begin{refsection}[citations.bib]
% \nocite{*}
% \printbibliography[heading=none]
% \end{refsection}

%----------------------------------------------------------------------------------------
%	SKILLS
%----------------------------------------------------------------------------------------
\section{Skills}
\begin{tabularx}{\linewidth}{@{}l X@{}}
Programming Languages & \normalsize{C/C++, Python, MATLAB}\\
Embedded Platforms & \normalsize{STM32 (H7/G4), NUC, Raspberry Pi}\\
Control Theory & \normalsize{PID, LQR, MPC, EKF, State Machine, Adaptive Control}\\
Communication Protocols & \normalsize{CAN, UART, SPI, I2C, DMA}\\
Development Tools & \normalsize{CMake, Git, Linux, ROS2, OpenCV, LaTeX}\\
\end{tabularx}

% \vfill
% \center{\footnotesize Last updated: \today}

\end{document}